% =====================================================================
% DRAFT PARAGRAPH FOR main.tex
% =====================================================================
%
% Purpose: establish that Koide (and by extension this paper's
% Soddy-Descartes construction) is part of a legitimate tradition of
% empirical mass formulas, several of which eventually received
% dynamical explanations. This is the "build credibility" paragraph
% the user requested.
%
% Recommended placement: near the end of Section 1 ("The Koide formula
% and its convergence"), immediately before or after the convergence
% table, or as a new subsection titled "Historical context" or
% "Tradition of empirical mass formulas."
%
% Alternative placement: in Section 8 ("Discussion"), as a paragraph
% titled "Historical precedent."
%
% The paragraph is ~250 words. It cites six references that should be
% added to the bibliography: Koide 1983, Rivero-Gsponer 2005,
% Sumino 2009, Gell-Mann 1961 (or CTSL-20 report), Okubo 1962,
% Barut 1979, Harari-Haut-Weyers 1978.
%
% =====================================================================

\paragraph{Historical context.}
The Koide relation sits within a tradition of empirical mass formulas
in particle physics. The Gell-Mann--Okubo sum rule for hadron
multiplets~\cite{GellMann1961,Okubo1962}, the Balmer series for
hydrogen spectral lines, and the Cabibbo angle~\cite{Cabibbo1963}
were each quantitative patterns discovered without dynamical
justification; in each case the empirical discovery preceded and
motivated the theoretical framework (SU(3) flavor symmetry,
quantum mechanics, the CKM matrix) that eventually explained it.
The Koide formula has a substantial literature spanning four
decades. Precursors include the Harari-Haut-Weyers quark mass
ansatz~\cite{HarariHautWeyers1978} and Barut's lepton mass
formula~\cite{Barut1979}, both of which predate Koide and also
relate charged fermion masses via simple algebraic expressions.
Geometric interpretations of Koide's $Q = 2/3$ have been developed
by Foot~\cite{Foot1994} (as the square cosine of the angle between
$(\sqrt{m_e}, \sqrt{m_\mu}, \sqrt{m_\tau})$ and the unit diagonal)
and by Kocik~\cite{Kocik2012} (who observed the formal similarity
to a generalized Descartes circle formula that the present paper
develops in detail). Extensions of the relation to
quarks~\cite{Kartavtsev2011,LiMa2016} and
neutrinos~\cite{RodejohannZhang2011,Cao2012}, dynamical
explanations via family gauge symmetry~\cite{Sumino2009a,Sumino2009b},
RG-running analyses~\cite{XingZhang2006,LiMa2016}, and
historical reviews~\cite{RiveroGsponer2005} have appeared.  The
persistence of the Koide prediction across a factor of 35
improvement in tau-mass measurement precision (Table~1) is the
signature that motivates treating the relation as reflecting
structure rather than coincidence; the Monte Carlo null test in
the companion paper~\cite{BrilliantMC2026} quantifies this
against a scale-invariant random-mass prior.

% =====================================================================
% BIBLIOGRAPHY ENTRIES TO ADD
% =====================================================================
%
% Add these entries to the \begin{thebibliography} block in main.tex.
% All are verified or marked for verification.
%
% \bibitem{GellMann1961}
%   M.~Gell-Mann, Caltech Synchrotron Laboratory Report CTSL-20 (1961).
%
% \bibitem{Okubo1962}
%   S.~Okubo, Prog.\ Theor.\ Phys.\ \textbf{27}, 949 (1962).
%
% \bibitem{Cabibbo1963}
%   N.~Cabibbo, Phys.\ Rev.\ Lett.\ \textbf{10}, 531 (1963).
%
% \bibitem{HarariHautWeyers1978}
%   H.~Harari, H.~Haut, J.~Weyers, Phys.\ Lett.\ B \textbf{78}, 459 (1978).
%
% \bibitem{Barut1979}
%   A.~O.~Barut, Phys.\ Rev.\ Lett.\ \textbf{42}, 1251 (1979).
%
% \bibitem{Foot1994}
%   R.~Foot, arXiv:hep-ph/9402242 (1994).
%
% \bibitem{RiveroGsponer2005}
%   A.~Rivero, A.~Gsponer, arXiv:hep-ph/0505220 (2005).
%
% \bibitem{XingZhang2006}
%   Z.-Z.~Xing, H.~Zhang, Phys.\ Lett.\ B \textbf{635}, 107 (2006).
%
% \bibitem{Sumino2009a}
%   Y.~Sumino, JHEP \textbf{0905}, 075 (2009), arXiv:0812.2103.
%
% \bibitem{Sumino2009b}
%   Y.~Sumino, Phys.\ Lett.\ B \textbf{671}, 477 (2009), arXiv:0812.2090.
%
% \bibitem{Kartavtsev2011}
%   A.~Kartavtsev, arXiv:1111.0480 (2011).
%
% \bibitem{RodejohannZhang2011}
%   W.~Rodejohann, H.~Zhang, Phys.\ Lett.\ B \textbf{698}, 152 (2011).
%
% \bibitem{Cao2012}
%   F.~G.~Cao, Phys.\ Rev.\ D \textbf{85}, 113003 (2012).
%
% \bibitem{Kocik2012}
%   J.~Kocik, arXiv:1201.2067 (2012).
%   % NOTE: arXiv number to verify before submission.
%
% \bibitem{LiMa2016}
%   N.~Li, B.-Q.~Ma et al., Eur.\ Phys.\ J.\ C \textbf{76}, 167 (2016),
%   arXiv:1601.00805.
%
% \bibitem{BrilliantMC2026}
%   A.~M.~Brilliant, companion Monte Carlo analysis, 2026. In preparation.
%
% =====================================================================
% END OF DRAFT
% =====================================================================
